\documentclass[../DoAn.tex]{subfiles}
\begin{document}

% Lưu ý: \textbf{Trước khi viết ĐATN, sinh viên cần đọc kỹ hướng dẫn và quy định chi tiết} về cách viết ĐATN trong Phụ lục A. Sinh viên tuân theo mẫu tài liệu này để viết báo cáo đồ án tốt nghiệp, vì tài liệu này đã được căn chỉnh, chỉnh sửa theo đúng chuẩn báo cáo kỹ thuật đồ án tốt nghiệp (ISO 7144:1986). Sinh viên viết trực tiếp vào file này, chỉ chỉnh sửa nội dung, và không viết trên file mới.

% \textbf{Khi đóng quyển ĐATN}, sinh viên cần lưu ý tuân thủ hướng dẫn ở phụ lục A.9

% \textbf{SV cần đặc biệt lưu ý cách hành văn}. Mỗi đoạn văn không được quá dài và cần có ý tứ rõ ràng, bao gồm duy nhất một ý chính và các ý phân tích bổ trợ để làm rõ hơn ý chính. Các câu văn trong đoạn phải đầy đủ chủ ngữ vị ngữ, cùng hướng đến chủ đề chung. Câu sau phải liên kết với câu trước, đoạn sau liên kết với đoạn trước. Trong văn phong khoa học, sinh viên không được dùng từ trong văn nói, không dùng các từ phóng đại, thái quá, các từ thiếu khách quan, thiên về cảm xúc, về quan điểm cá nhân như “tuyệt vời”, “cực hay”, “cực kỳ hữu ích”, v.v. Các câu văn cần được tối ưu hóa, đảm bảo rất khó để thể thêm hoặc bớt đi được dù chỉ một từ. Cách diễn đạt cần ngắn gọn, súc tích, không dài dòng.

% Mẫu ĐATN này được thiết kế phù hợp nhất với đa số các đề tài xây dựng phần mềm ứng dụng. Với các dạng đề tài khác (giải pháp, nghiên cứu, phần mềm đặc thù, v.v.), sinh viên dựa trên cấu trúc và hướng dẫn của báo cáo này để đề xuất và trao đổi với giáo viên hướng dẫn để thiết kế khung báo cáo đồ án cho phù hợp. Sinh viên lưu ý \textbf{trong mọi trường hợp, SV luôn phải sử dụng định dạng báo cáo này, và phải đọc kỹ toàn bộ các hướng dẫn từ đầu tới cuối.} Các hướng dẫn không chỉ áp dụng riêng cho đề tài ứng dụng, mà còn phù hợp với các dạng đề tài khác. Ngoài ra, trong mẫu ĐATN này đã được tích hợp một số hướng dẫn dành riêng cho đề tài nghiên cứu.

% Chương 1 có độ dài từ 3 đến 6 trang với các nội dung sau đây

\section{Đặt vấn đề}
\label{section:1.1}
% Khi đặt vấn đề, sinh viên cần làm nổi bật mức độ cấp thiết, tầm quan trọng và/hoặc quy mô của bài toán của mình.

% Gợi ý cách trình bày cho sinh viên: Xuất phát từ tình hình thực tế gì, dẫn đến vấn đề hoặc bài toán gì. Vấn đề hoặc bài toán đó, nếu được giải quyết, đem lại lợi ích gì, cho những ai, còn có thể được áp dụng vào các lĩnh vực khác nữa không. Sinh viên cần lưu ý phần này chỉ trình bày vấn đề, tuyệt đối không trình bày giải pháp.

Trong mối cảnh chuyển đổi số ngày càng phát triển, các cơ sở giáo dục tiểu học cũng đang đối mặt với nhu cầu cao về tự động hóa các quy trình, thủ tục hành chính để tăng hiệu suất công việc, trong đó có quy trình sắp xếp thời khoá biểu. Thời khóa biểu là cơ sở để nhà trường tổ chức hoạt động giảng dạy, phân công giáo viên và sử dụng cơ sở vật chất trong từng tuần. Để xây dựng một thời khóa biểu, người phụ trách không chỉ cần phân bổ các môn học vào từng tiết mà còn phải đồng thời xem xét số tiết yêu cầu của mỗi môn, giáo viên được phân công, lịch học của từng lớp và khả năng sử dụng các phòng chức năng.

Tuy nhiên, hiện nay tại nhiều trường tiểu học, việc sắp xếp thời khoá biểu vẫn được thực hiện thủ công bằng bảng tính dựa trên kinh nghiệm cá nhân của người phụ trách. Mặc dù phương pháp này cho phép người lập lịch chủ động sắp xếp và điều chỉnh từng tiết học, quá trình thực hiện vẫn tồn tại nhiều hạn chế như khó kiểm soát đồng thời lịch dạy của giáo viên, lịch học của lớp, số tiết yêu cầu của từng môn và khả năng sử dụng phòng chức năng. Người phụ trách phải đối chiếu nhiều nguồn dữ liệu bằng tay nên mất nhiều thời gian và dễ bỏ sót các xung đột. Khi thông tin giáo viên hoặc phân công giảng dạy thay đổi, nhiều vị trí trong thời khóa biểu phải được rà soát và điều chỉnh lại, làm tăng khối lượng công việc và có thể phát sinh thêm các xung đột liên quan. Đặc biệt với các trường tiểu học quy mô nhỏ, việc đầu tư vào các phần mềm xếp thời khóa biểu thương mại còn gặp hạn chế về chi phí.

Từ thực tế đó, đề tài "EduSchedule — Hệ thống hỗ trợ xếp thời khóa biểu cho trường tiểu học" được thực hiện nhằm xây dựng một ứng dụng web miễn phí, cho phép ban giám hiệu trường tiểu học xếp thời khóa biểu theo hai phương thức: thủ công trên giao diện lưới trực quan, hoặc tự động bằng thuật toán với kết quả do người dùng xem xét tính hợp lý và quyết định áp dụng. Hệ thống đồng thời phát hiện xung đột theo thời gian thực, giúp ban giám hiệu vận hành công tác xếp thời khóa biểu một cách chính xác và tiết kiệm công sức hơn so với phương pháp thủ công truyền thống.

\section{Mục tiêu và phạm vi đề tài}
\label{section:1.2}
% Sinh viên trước tiên cần trình bày tổng quan các kết quả của các nghiên cứu hiện nay cho bài toán giới thiệu ở phần \ref{section:1.1} (đối với đề tài nghiên cứu), hoặc về các sản phẩm hiện tại/về nhu cầu của người dùng (đối với đề tài ứng dụng). Tiếp đến, sinh viên tiến hành so sánh và đánh giá tổng quan các sản phẩm/nghiên cứu này.

% Dựa trên các phân tích và đánh giá ở trên, sinh viên khái quát lại các hạn chế hiện tại đang gặp phải. Trên cơ sở đó, sinh viên sẽ hướng tới giải quyết vấn đề cụ thể gì, khắc phục hạn chế gì, phát triển phần mềm \textbf{có các chức năng chính gì}, tạo nên đột phá gì, v.v.

% Trong phần này, sinh viên lưu ý chỉ trình bày tổng quan, không đi vào chi tiết của vấn đề hoặc giải pháp. Nội dung chi tiết sẽ được trình bày trong các chương tiếp theo, đặc biệt là trong Chương 5.
\subsection{Mục tiêu}
\label{section:1.2.1}

Mục tiêu tổng quát của đề tài là xây dựng một hệ thống web hỗ trợ lập thời khóa biểu cho trường tiểu học, đáp ứng các yêu cầu nghiệp vụ liên quan tới việc sắp xếp thời khoá biểu của ban giám hiệu và nhu cầu tra cứu thời khóa biểu của giáo viên.
Để đạt được mục tiêu tổng quát trên, đề tài tập trung thực hiện các mục tiêu cụ thể sau:
\begin{itemize}
    \item \textbf{Nhóm chức năng quản lý và phân công giảng dạy:}
    Xây dựng các chức năng quản lý năm học, giáo viên, lớp học, môn học và phòng chức năng; đồng thời hỗ trợ người phụ trách phân công giáo viên giảng dạy các môn học theo từng lớp.

    \item \textbf{Nhóm chức năng lập thời khóa biểu:}
    Xây dựng giao diện thời khóa biểu dạng lưới, hỗ trợ người phụ trách sắp xếp và điều chỉnh tiết học thủ công, kiểm tra các xung đột liên quan đến giáo viên, lớp học và phòng chức năng, đồng thời hỗ trợ tự động phân bổ các tiết còn thiếu vào những vị trí trống mà không làm thay đổi các tiết đã được sắp xếp trước.

    \item \textbf{Nhóm chức năng tra cứu và xuất thời khóa biểu:}
    Cho phép xem thời khóa biểu theo lớp, giáo viên hoặc khối; hỗ trợ giáo viên tra cứu thời khóa biểu được nhà trường công bố và cho phép người phụ trách xuất thời khóa biểu ra file Excel.

    \item \textbf{Chất lượng và an toàn hệ thống:}
    Xây dựng hệ thống web có giao diện trực quan, tổ chức mã nguồn theo kiến trúc client--server, triển khai cơ chế xác thực để bảo vệ các chức năng quản lý, đồng thời xây dựng các kịch bản kiểm thử cho những nghiệp vụ chính của hệ thống.
\end{itemize}

\subsection{Phạm vi đề tài}
\label{section:1.2.2}

Đề tài được thực hiện trong phạm vi hỗ trợ công tác lập thời khóa biểu cho một trường tiểu học, bao gồm các lớp từ khối 1 đến khối 5. Thời khóa biểu được quản lý theo từng năm học, kì học và từng tuần, với các thông tin liên quan đến giáo viên, lớp học, môn học, phòng chức năng và phân công giảng dạy.

Hệ thống phục vụ hai nhóm người dùng gồm người phụ trách lập thời khóa biểu và người dùng vãng lai. Người phụ trách có quyền quản lý thông tin và thực hiện các nghiệp vụ liên quan đến lập thời khóa biểu. Người dùng vãng lai, chủ yếu là giáo viên, chỉ có quyền tra cứu thời khóa biểu được nhà trường công bố và không được phép thay đổi dữ liệu của hệ thống.

Trong phạm vi xếp lịch, hệ thống hỗ trợ người dùng kết hợp giữa sắp xếp thủ công và xếp tự động. Các tiết đã được sắp xếp trước được giữ cố định, còn chức năng xếp tự động chỉ xử lý những tiết còn thiếu. Kết quả tự động tạo ra cần được người phụ trách kiểm tra và xác nhận trước khi lưu chính thức.

Đề tài chưa hỗ trợ tài khoản riêng cho giáo viên, học sinh và phụ huynh; gửi thông báo, nhắc lịch hoặc đồng bộ với các ứng dụng lịch bên ngoài. Hệ thống chỉ áp dụng cho trường tiểu học và không bảo đảm luôn tìm được phương án xếp đầy đủ khi các ràng buộc đầu vào không thể được thỏa mãn đồng thời.

\section{Định hướng giải pháp}
\label{section:1.3}
% Từ việc xác định rõ nhiệm vụ cần giải quyết ở phần \ref{section:1.2}, sinh viên đề xuất định hướng giải pháp của mình theo trình tự sau: (i) Sinh viên trước tiên trình bày sẽ giải quyết vấn đề theo định hướng, phương pháp, thuật toán, kỹ thuật, hay công nghệ nào; Tiếp theo, (ii) sinh viên mô tả ngắn gọn giải pháp của mình là gì (khi đi theo định hướng/phương pháp nêu trên); và sau cùng, (iii) sinh viên trình bày đóng góp chính của đồ án là gì, kết quả đạt được là gì.

% Sinh viên lưu ý không giải thích hoặc phân tích chi tiết công nghệ/thuật toán trong phần này. Sinh viên chỉ cần nêu tên định hướng công nghệ/thuật toán, mô tả ngắn gọn trong một đến hai câu và giải thích nhanh lý do lựa chọn.

Trên cơ sở các mục tiêu đã xác định, đề tài lựa chọn định hướng xây dựng một ứng dụng web theo mô hình kiến trúc client--server, tách biệt giữa tầng giao diện người dùng và tầng xử lý nghiệp vụ phía máy chủ, giao tiếp thông qua các API theo phong cách RESTful. Định hướng này giúp hệ thống dễ mở rộng, bảo trì và phù hợp với yêu cầu quản lý, lập thời khóa biểu trên nền tảng web.

\begin{itemize}
    \item \textbf{Về phía máy chủ (backend):}
    Sử dụng Java 21 cùng framework Spring Boot để xây dựng các API và xử lý nghiệp vụ quản lý năm học, giáo viên, lớp học, môn học, phòng chức năng, phân công giảng dạy và thời khóa biểu. Spring Data JPA được sử dụng để tương tác với cơ sở dữ liệu PostgreSQL, trong khi Spring Security kết hợp với JWT đảm nhiệm việc xác thực và bảo vệ các chức năng quản lý.

    \item \textbf{Về phía giao diện (frontend):}
    Sử dụng Next.js, TypeScript và Tailwind CSS để xây dựng giao diện tương tác. Hệ thống cung cấp giao diện thời khóa biểu dạng lưới, cho phép người phụ trách xem, sắp xếp và điều chỉnh tiết học, đồng thời hỗ trợ hiển thị thời khóa biểu theo lớp, giáo viên hoặc khối.

    \item \textbf{Về chức năng xếp thời khóa biểu tự động:}
    Sử dụng phương pháp khởi tạo lịch kết hợp với Timefold Solver để phân bổ các tiết học còn thiếu vào những vị trí trống theo các ràng buộc về giáo viên, lớp học và phòng chức năng. Các tiết đã được người dùng sắp xếp trước được giữ cố định trong quá trình xếp tự động. Kết quả được hiển thị để người phụ trách kiểm tra, điều chỉnh và xác nhận trước khi lưu.

    \item \textbf{Về tra cứu và xuất thời khóa biểu:}
    Hệ thống hỗ trợ giáo viên tra cứu thời khóa biểu được nhà trường công bố với tư cách người dùng vãng lai. Bên cạnh đó, thư viện SheetJS được sử dụng để xuất thời khóa biểu theo lớp, giáo viên hoặc khối ra file Excel.
\end{itemize}

\section{Bố cục đồ án}
\label{section:1.4}
% Sinh viên cần viết mô tả thành đoạn văn đầy đủ về nội dung chương. Tuyệt đối không viết ý hay gạch đầu dòng. Chương 1 không cần mô tả trong phần này.
%
% Ví dụ tham khảo mô tả chương trong phần bố cục đồ án tốt nghiệp: Chương *** trình bày đóng góp chính của đồ án, đó là một nền tảng ABC cho phép khai phá và tích hợp nhiều nguồn dữ liệu, trong đó mỗi nguồn dữ liệu lại có định dạng đặc thù riêng. Nền tảng ABC được phát triển dựa trên khái niệm DEF, là các module ngữ nghĩa trợ giúp người dùng tìm kiếm, tích hợp và hiển thị trực quan dữ liệu theo mô hình cộng tác và mô hình phân tán.
%
% Chú ý: Trong phần nội dung chính, mỗi chương của đồ án nên có phần Tổng quan và Kết chương. Hai phần này đều có định dạng văn bản “Normal”, sinh viên không cần tạo định dạng riêng, ví dụ như không in đậm/in nghiêng, không đóng khung, v.v.
%
% Trong phần Tổng quan của chương N, sinh viên nên có sự liên kết với chương N-1 rồi trình bày sơ qua lý do có mặt của chương N và sự cần thiết của chương này trong đồ án. Sau đó giới thiệu những vấn đề sẽ trình bày trong chương này là gì, trong các đề mục lớn nào.
%
% Ví dụ về phần Tổng quan: Chương 3 đã thảo luận về nguồn gốc ra đời, cơ sở lý thuyết và các nhiệm vụ chính của bài toán tích hợp dữ liệu. Chương 4 này sẽ trình bày chi tiết các công cụ tích hợp dữ liệu theo hướng tiếp cận “mashup”. Với mục đích và phạm vi của đề tài, sáu nhóm công cụ tích hợp dữ liệu chính được trình bày bao gồm: (i) nhóm công cụ ABC trong phần 4.1, (ii) nhóm công cụ DEF trong phần 4.2, nhóm công cụ GHK trong phần 4.3, v.v.
%
% Trong phần Kết chương, sinh viên đưa ra một số kết luận quan trọng của chương. Những vấn đề mở ra trong Tổng quan cần được tóm tắt lại nội dung và cách giải quyết/thực hiện như thế nào. Sinh viên lưu ý không viết Kết chương giống hệt Tổng quan. Sau khi đọc phần Kết chương, người đọc sẽ nắm được sơ bộ nội dung và giải pháp cho các vấn đề đã trình bày trong chương. Trong Kết chương, Sinh viên nên có thêm câu liên kết tới chương tiếp theo.
%
% Ví dụ về phần Kết chương: Chương này đã phân tích chi tiết sáu nhóm công cụ tích hợp dữ liệu. Nhóm công cụ ABC và DEF thích hợp với những bài toán tích hợp dữ liệu phạm vi nhỏ. Trong khi đó, nhóm công cụ GHK lại chứng tỏ thế mạnh của mình với những bài toán cần độ chính xác cao, v.v. Từ kết quả nghiên cứu và phân tích về sáu nhóm công cụ tích hợp dữ liệu này, tôi đã thực hiện phát triển phần mềm tự động bóc tách và tích hợp dữ liệu sử dụng nhóm công cụ GHK. Phần này được trình bày trong chương tiếp theo – Chương 5.

Phần còn lại của báo cáo đồ án tốt nghiệp này được tổ chức như sau.

Chương 2 trình bày quá trình khảo sát và phân tích yêu cầu hệ thống. Dựa trên thực trạng xếp thời khóa biểu tại các trường tiểu học công lập, chương này xác định các tác nhân, phân tích use case và đặc tả chi tiết các chức năng cốt lõi của hệ thống, đồng thời làm rõ các yêu cầu phi chức năng cần đáp ứng.

Chương 3 giới thiệu các công nghệ được sử dụng để xây dựng hệ thống, bao gồm Next.js và React phía frontend, Spring Boot và Spring Security phía backend, PostgreSQL làm hệ quản trị cơ sở dữ liệu quan hệ, cùng các thư viện hỗ trợ xác thực và xuất dữ liệu. Chương này phân tích lý do lựa chọn từng công nghệ trong bối cảnh yêu cầu của đề tài.

Chương 4 trình bày quá trình thiết kế, xây dựng và đánh giá hệ thống. Nội dung bao gồm thiết kế kiến trúc tổng thể, thiết kế giao diện, thiết kế lớp và cơ sở dữ liệu; quá trình hiện thực các chức năng; giải pháp xếp thời khóa biểu thủ công và tự động; kiểm thử chức năng; cùng phương thức triển khai sản phẩm.

Chương 5 tập trung vào các giải pháp và đóng góp nổi bật của đồ án, bao gồm tính năng xếp thời khóa biểu tự động kết hợp điều chỉnh thủ công trên giao diện lưới tương tác, cơ chế phát hiện xung đột tiết học theo thời gian thực, và cách tiếp cận triển khai hệ thống không phát sinh chi phí cho nhà trường.

Chương 6 đưa ra kết luận về mức độ hoàn thành các mục tiêu đề ra, đánh giá những kết quả đạt được và hạn chế còn tồn tại, đồng thời đề xuất các hướng phát triển tiếp theo nhằm nâng cao tính hoàn thiện và khả năng ứng dụng thực tế của hệ thống.

\end{document}